\chapter*{Введение}
\addcontentsline{toc}{chapter}{Введение}
\label{ch:intro}

\textbf{Figma} — облачный графический редактор интерфейсов, ставший стандартом
в индустрии UI/UX-дизайна благодаря векторной природе, совместной работе в
реальном времени и широкой экосистеме плагинов. После знакомства с базовыми
инструментами Figma в лабораторной работе~10 на примере одностраничного
лендинга, следующим естественным шагом является разработка многостраничного
сайта с единым визуальным языком и адаптивной вёрсткой каждой страницы.

В качестве темы выбран сайт чемпионата по фиджитал-спорту: новому виду
соревнований, в которых классическая физическая дисциплина дополняется
её цифровым (видеоигровым) аналогом. В работе рассматривается
гипотетический турнир \textbf{PHYGITAL~FC} по киберфутболу, проводимый
в Новосибирске в 2024~году. Сайт должен решать три ключевые задачи:
информационное сопровождение спортивной части (команды, участники,
турнирная таблица), общее освещение события (новости, медиа, расписание)
и коммерческую поддержку (продажа билетов и сувенирной продукции).

Цель работы — освоить проектирование многостраничного веб-интерфейса в
Figma с соблюдением единого фирменного стиля и адаптивной вёрсткой каждой
страницы для трёх типовых ширин экрана.

Задачи:
\begin{enumerate}
  \item Разработать дизайн-концепцию главной страницы сайта чемпионата.
  \item Разработать страницу команды-участницы со списком игроков.
  \item Разработать страницу карточки участника (профиль игрока).
  \item Разработать страницу турнирной таблицы (групповой этап и
        плей-офф).
  \item Разработать страницу магазина с двумя секциями: продажа билетов
        (с выбором типа: только просмотр или просмотр + дегустация) и
        каталог сувенирной продукции с символикой чемпионата.
  \item Настроить колоночную \texttt{Layout grid} и привязки
        \texttt{Constraints} для всех ключевых блоков каждой страницы.
  \item Подготовить адаптивные версии каждой из пяти страниц для двух
        дополнительных ширин: планшет (\texttt{768\,px}) и мобильный
        телефон (\texttt{360\,px}).
\end{enumerate}

\textbf{Стенд:}
\begin{itemize}
  \item Веб-приложение Figma (\texttt{figma.com}), доступ через
        современный браузер; альтернатива — desktop-приложение Figma
        для Windows/macOS.
  \item Плагин \texttt{html.to.design} для импорта подготовленного
        HTML-каркаса макета во фреймы Figma.
  \item Плагин \texttt{Unsplash} для подбора тематических фотографий
        (стадион, болельщики, портреты игроков, мерч).
  \item Шрифты Google Fonts: \texttt{Space~Grotesk} (заголовки),
        \texttt{Inter} (текст интерфейса), \texttt{JetBrains~Mono}
        (числовая статистика).
\end{itemize}

В отличие от лабораторных работ~4--9, стенд не использует виртуальные
машины и не требует bash-автоматизации: вся работа ведётся в графическом
интерфейсе Figma. Поэтому в каталоге \texttt{lab11/} отсутствуют файлы
\texttt{config.sh} и сценарии настройки сервисов.

\endinput
